% Diagrama: Figura B.2 - Componentes del backend
% TikZ vectorial. Se incluye con \input desde apendices.tex.
% Una sola aplicación (monolito modular): la API REST (controladores) recibe
% las peticiones y las pasa a los módulos expuestos (fila superior); los
% módulos internos (fila inferior) los usan. La simulación se ejecuta en un
% proceso aparte (worker) de la misma aplicación y reutiliza el matching y el
% Gemelo Digital. Todos los módulos acceden a la base de datos a través de la
% capa de acceso a datos (base del diagrama), sin flechas para no saturarlo.

\begin{tikzpicture}[x=1cm,y=1cm]
  \useasboundingbox (0.0,1.55) rectangle (15.55,8.35);
  \tikzset{bus/.style={draw=Color2!75, line width=0.45pt, dashed}}
  \def\ya{5.35}  % módulos expuestos por la API
  \def\yb{3.85}  % módulos internos y worker

  % ---------------- frontera: una sola aplicación ----------------
  \draw[draw=Color2!55, dashed, rounded corners=4pt, fill=Color2!2]
    (0.40,1.75) rectangle (12.70,7.55);
  \node[font=\scriptsize\bfseries, text=Color2, anchor=south west] at (0.55,7.60)
    {\textit{Backend}: una aplicación Node.js + NestJS};

  % ---------------- componentes ----------------
  \node[componente, minimum width=4.40cm] (api) at (6.52,6.85) {API REST (controladores)};

  \foreach \n/\t/\x in {usu/Usuarios/1.36, via/Viajes/3.08, sol/Solicitudes/4.80,
      inc/Incidencias/6.52, seg/Seguimiento/8.24, mat/\textit{Matching}/9.96, lug/Lugares/11.68}
    \node[componente, minimum width=1.60cm] (\n) at (\x,\ya) {\t};

  \node[componente, minimum width=1.60cm] (not) at (4.80,\yb) {Notificaciones};
  \node[componente, minimum width=1.60cm] (dt)  at (6.95,\yb) {Gemelo Digital};
  \node[componente, minimum width=1.60cm] (met) at (8.85,\yb) {Métricas};

  % worker de simulación: proceso aparte de la misma aplicación
  \draw[draw=Color2!55, dashed, rounded corners=3pt] (0.52,3.30) rectangle (2.20,4.55);
  \node[font=\tiny, text=Color2, anchor=north west] at (0.56,4.53) {\textit{worker}};
  \node[componente, minimum width=1.50cm] (sim) at (1.36,3.80) {Simulación};

  \node[componente, minimum width=12.00cm] (repo) at (6.55,2.35) {Acceso a datos (TypeORM) y registro de eventos};

  % ---------------- elementos externos ----------------
  \node[componente, minimum width=2.40cm, fill=white] (fe)   at (14.30,6.85) {\textit{Frontend} y panel};
  \node[componente, minimum width=2.40cm, fill=white] (cc)   at (14.30,\ya) {CartoCiudad (IGN)};
  \node[componente, minimum width=2.40cm, fill=white] (osrm) at (14.30,\yb) {Motor de rutas\\ (OSRM)};
  \node[componente, minimum width=2.40cm, fill=white] (bd)   at (14.30,2.35) {PostgreSQL + PostGIS};

  % ---------------- dependencias ----------------
  \draw[dep] (fe.west) -- (api.east);
  % API REST -> módulos expuestos
  \draw[bus] (api.south) -- (6.52,6.15);
  \draw[bus] (1.36,6.15) -- (11.68,6.15);
  \foreach \c in {usu,via,sol,inc,seg,mat,lug} \draw[dep] (\c.north |- 0,6.15) -- (\c.north);

  % módulos -> módulos internos y servicios externos
  \draw[dep] (sol.south) -- (not.north);
  \draw[dep] (lug.east) -- (cc.west);
  \draw[dep] (mat.south) -- (9.96,\yb) -- (osrm.west);
  % el worker reutiliza el matching y el Gemelo Digital
  \draw[dep] (sim.south) -- (1.36,3.12) -- (9.70,3.12) -- (9.70,0 |- mat.south);
  \draw[dep] (6.95,3.12) -- (dt.south);

  % persistencia
  \draw[dep] (repo.east) -- (bd.west);
\end{tikzpicture}
